\begin{comment}
\begin{frame}{Epílogo: Un desafío estructural}
\begin{itemize}
    \item La crisis fiscal actual es estructural: incapacidad de aumentar recaudo y retornar a la senda de gasto pre-pandemia.
    \item Requiere solución de Estado: superar barreras políticas, legales y constitucionales.
    \item Urgencia: evitar que la deuda pública alcance niveles inmanejables y que se intensifiquen presiones de gasto (intereses, SGP, pensiones, salud, personal).
    \item Retrasar el ajuste aumenta el costo social, económico y político.
\end{itemize}
\end{frame}

\begin{frame}{Ruta para la consolidación fiscal}
\begin{itemize}
    \item Ajuste superior a 4\% del PIB: reformas tributarias sucesivas, racionalización del gasto en salud, personal y pensiones.
    \item Compatibilidad de la Ley de Competencias del SGP con sostenibilidad fiscal.
    \item Evitar que nuevos ingresos tributarios se transfieran automáticamente vía SGP.
    \item Revisión profunda de beneficios tributarios, especialmente en IVA, con mecanismos de compensación a hogares vulnerables.
\end{itemize}
\end{frame}
\end{comment}

\begin{frame}{Comentarios Finales}
\begin{itemize}
%    \item Reformas institucionales.
    %\item Crecimiento económico sostenido como condición para estabilidad fiscal.
    \item Agenda pro-crecimiento: 
    \begin{itemize}
        \item Libre competencia
        \item Flexibilización laboral
        \item Modernización financiera 
        \item Asignación del riesgo
        \item Incentivos a la innovación
    \end{itemize}
    \item ¿Disyuntiva entre estabilidad financiera y control de la inflación?
    \item ¿Cómo vender la idea de pagar más impuestos y reducir el gasto?
\end{itemize}
\end{frame}
